\documentclass[12pt,a4paper]{article}
\usepackage[utf8]{inputenc}
\usepackage[T1]{fontenc}
\usepackage{geometry}
\usepackage{setspace}
\usepackage{tabularx}
\usepackage{booktabs}
\usepackage{array}
\geometry{margin=1in}
\setstretch{1.25}
\setlength{\parindent}{2em}
\setlength{\parskip}{0.4em}

\begin{document}

\begin{titlepage}
\vspace*{\fill}
\begin{center}
{\LARGE \textbf{Bachelor Project Plan}\par}
\end{center}
\vspace{2em}
\raggedright
{\large \textbf{Name:} Xuanlin Chen\par}
{\large \textbf{Advisor:} Prof.\ Piotr Didyk\par}
{\large \textbf{Title:} Fluid Simulation and Real-Time Rendering of Shallow Water Equations\par}
\vspace*{\fill}
\end{titlepage}

\section*{Project Objectives}
This project aims to simulate water-surface behavior by numerically solving the shallow water equations (SWE). It focuses on dynamic fluctuations, user interaction, and real-time rendering for large-scale water scenes while preserving physical plausibility.

\section*{Work Plan}
\begin{table}[h]
\centering
\renewcommand{\arraystretch}{1.2}
\begin{tabularx}{\textwidth}{>{\centering\arraybackslash}p{1cm} X X}
\toprule
\textbf{Week} & \textbf{Planned Activity} & \textbf{Expected Result} \\
\midrule
W1 & Setup environment and basic shaders & The screen can render a static plane with simple colors. \\
W2 & Implement SWE base update cycle & The fundamental changes in the height field can be observed. \\
W3 & Staggered grid \& advection implementation & Stable fluid motion \\
W4 & CFL condition \& fluctuation handling & Numerically stable simulation \\
W5 & Blinn-Phong shading implementation & Giving the waveform a sense of depth and light and shadow feedback. \\
W6 & Real-time vertex height displacement & Dynamic water mesh \\
W7 & Implement Boundary conditions & Constrained fluid behavior \\
W8 & Environment mapping \& Fresnel effect & Enhanced water aesthetics \\
W9 & Solid-fluid coupling & Interaction with dynamic objects \\
W10 & Particle system / Foam map & Visual detail enrichment \\
W11 & Performance optimization & Real-time frame rate stability \\
W12 & Final documentation & Completed Thesis Report \\
\bottomrule
\end{tabularx}
\end{table}
\end{document}